% from ega3-1-fr.tex
\oldpage[III]{131}

Soit $Y'$ le $Y$-schéma fini au-dessus de $Y$ tel que
$\mathcal A(Y')=\mathcal A(X)$, qui est déterminé à un $Y$-isomorphisme
près (II, 1.3.1 et 6.1.3)~; si $f'=\mathcal A(e)$ est le $Y$-morphisme
$X\to Y'$ déduit de l'isomorphisme identique
$e:\mathcal A(Y')\to\mathcal A(X)$ (II, 1.2.7), $f'$ est propre,
$f'_*(\mathcal O_X)$ est isomorphe à $\mathcal O_{Y'}$, et les fibres
$f'^{-1}(y')$ du morphisme $f'$ sont connexes et non vides pour tout
$y'\in Y'$.
\end{theorem}

\begin{proof}
Soit $g:Y'\to Y$ le morphisme structural. Pour prouver que l'homomorphisme
$\theta:\mathcal O_{Y'}\to f'_*(\mathcal O_X)$ entrant dans la définition
du morphisme $f'$ est bijectif, il suffit, puisque $Y'$ est affine sur $Y$,
de prouver que
\[
g_*(\theta):g_*(\mathcal O_{Y'})\longrightarrow
g_*(f'_*(\mathcal O_X))=f_*(\mathcal O_X)
\]
est l'identité (II, 1.4.2)~; mais cela résulte des définitions puisque
$g_*(\mathcal O_{Y'})=\mathcal A(Y')$ et
$f_*(\mathcal O_X)=\mathcal A(X)$. Le fait que $\mathcal A(X)$ est
cohérent est un cas particulier du th. de finitude
(\hyperref[III.3.2.1-fr]{3.2.1}). Comme $f$ est propre et $g$ séparé,
$f'$ est propre (II, 5.4.3, (i))~; pour terminer la démonstration de
(\hyperref[III.4.3.1-fr]{4.3.1}), il suffit donc d'en prouver le
\end{proof}

\begin{corollary}[4.3.2]
\label{III.4.3.2-fr}
Sous les hypothèses de (\hyperref[III.4.3.1-fr]{4.3.1}), supposons de plus
que $f_*(\mathcal O_X)$ soit isomorphe à $\mathcal O_Y$. Alors les fibres
$f^{-1}(y)$ de $f$ sont connexes et non vides pour tout $y\in Y$.
\end{corollary}

\begin{proof}
L'hypothèse que $f_*(\mathcal O_X)$ est isomorphe à $\mathcal O_Y$
entraîne déjà que $f$ est dominant, donc surjectif puisque $f$ est une
application fermée. On peut se ramener, comme dans
(\hyperref[III.4.2.1-fr]{4.2.1}), au cas où $y$ est fermé dans $Y$~;
$f^{-1}(y)$ étant un espace noethérien, a un nombre fini de composantes
connexes, et c'est l'espace sous-jacent du complété $\widehat X$ le long de
$f^{-1}(y)$. Si $Z_i$ ($1\leq i\leq n$) sont ces composantes connexes, il
est clair que $\Gamma(\widehat X,\mathcal O_{\widehat X})$ est composé
direct des anneaux $\Gamma(Z_i,\mathcal O_{\widehat X})$, et chacun de ces
derniers n'est pas réduit à $0$, puisque la section unité est distincte de
$0$ en chaque point de $\widehat X$. Or, si on applique
(\hyperref[III.4.1.5-fr]{4.1.5}) à $\mathcal F=\mathcal O_X$, dont le
complété le long de $f^{-1}(y)$ est $\mathcal O_{\widehat X}$, on voit que
$\Gamma(\widehat X,\mathcal O_{\widehat X})$ est isomorphe au séparé
complété $\mathfrak m_y$-adique $\widehat{\mathcal O}_y$ de l'anneau local
$\mathcal O_y$~; c'est donc un anneau local qui ne peut être composé
direct de plusieurs anneaux non réduits à $0$ (sans quoi il aurait
plusieurs idéaux maximaux distincts). On a donc $n=1$, ce qui démontre le
corollaire.
\end{proof}

\begin{corollary}[4.3.3]
\label{III.4.3.3-fr}
Sous les hypothèses de (\hyperref[III.4.3.1-fr]{4.3.1}), pour tout $y\in Y$,
l'ensemble des composantes connexes de la fibre $f^{-1}(y)$ est en
correspondance biunivoque avec l'ensemble fini des points de la fibre
$g^{-1}(y)$, où $g:Y'\to Y$ est le morphisme structural (autrement dit,
l'ensemble des idéaux maximaux de $(f_*(\mathcal O_X))_y$).
\end{corollary}

\begin{proof}
Puisque $Y'$ est fini sur $Y$, on sait en effet que $g^{-1}(y)$ est un
espace fini discret (II, 6.1.7). Comme
$f^{-1}(y)=f'^{-1}(g^{-1}(y))$, le corollaire résulte de cette remarque et
de (\hyperref[III.4.3.1-fr]{4.3.1}).
\end{proof}

On a ainsi une interprétation remarquable du $Y$-préschéma $Y'$ défini
dans (\hyperref[III.4.3.1-fr]{4.3.1}). La factorisation $f=g\circ f'$ du
morphisme propre $f$ est analogue à la factorisation obtenue par K.~Stein
pour les applications holomorphes d'espaces analytiques, et nous
l'appellerons par la suite la \emph{factorisation de Stein} de $f$.

\begin{remark}[4.3.4]
\label{III.4.3.4-fr}
Soit $k$ une extension du corps $k(y)$~: si le préschéma
\[
f^{-1}(y)\otimes_{k(y)}k=X\times_Y\operatorname{Spec}(k)
\]
est connexe, il en est de même de $f^{-1}(y)$, qui en est l'image par un
morphisme de projection (I, 3.4.7). Nous dirons que, pour un morphisme
$f:X\to Y$ de préschémas et un point $y\in Y$, la fibre $f^{-1}(y)$ est
\emph{géométriquement connexe}, si pour toute extension $k$ de $k(y)$, le
préschéma
$f^{-1}(y)\otimes_{k(y)}k=X\times_Y\operatorname{Spec}(k)$ est connexe.

